

\bta{2009}

\begin{enumerate}
	%\renewcommand{\labelenumi}{\arabic{enumi}.}
	% A(\Alph) a(\alph) I(\Roman) i(\roman) 1(\arabic)
	%设定全局标号series=example	%引用全局变量resume=example
	%[topsep=-0.3em,parsep=-0.3em,itemsep=-0.3em,partopsep=-0.3em]
	%可使用leftmargin调整列表环境左边的空白长度 [leftmargin=0em]
	\item
（$30'$）已知在$(l^2,l_z)$的表象中，
$$l_x=\frac{\hbar}{\sqrt2}\begin{vmatrix}0&1&0\\1&0&1\\0&1&0\end{vmatrix}$$
求：
\begin{enumerate}
	%\renewcommand{\labelenumi}{\arabic{enumi}.}
	% A(\Alph) a(\alph) I(\Roman) i(\roman) 1(\arabic)
	%设定全局标号series=example	%引用全局变量resume=example
	%[topsep=-0.3em,parsep=-0.3em,itemsep=-0.3em,partopsep=-0.3em]
	%可使用leftmargin调整列表环境左边的空白长度 [leftmargin=0em]
	\item
$l_x$的本征值和相应的本征函数；

\item 
$l_y$的矩阵表示。
\end{enumerate}

\banswer{
	题给矩阵对应 $l=1$，基底顺序取 $|1,1\rangle,|1,0\rangle,|1,-1\rangle$。
	\begin{enumerate}
		\item 写 $L_x=\hbar M/\sqrt2$，则
		\[
		 \det(M-\eta I)=-\eta(\eta^2-2)=0.
		\]
		故 $L_x$ 的本征值为 $+\hbar,0,-\hbar$。对 $\eta=\pm\sqrt2$，
		方程 $Mv=\eta v$ 给出 $v_1=v_3$、$v_2=\eta v_1$；
		对 $\eta=0$，给出 $v_2=0$、$v_3=-v_1$。归一化后列矢为
		\[
			\chi_{+}=\begin{pmatrix}\frac12\\[2pt]\frac1{\sqrt2}\\[2pt]\frac12\end{pmatrix},
			\qquad
			\chi_{0}=\begin{pmatrix}\frac1{\sqrt2}\\[2pt]0\\[2pt]-\frac1{\sqrt2}\end{pmatrix},
			\qquad
			\chi_{-}=\begin{pmatrix}\frac12\\[2pt]-\frac1{\sqrt2}\\[2pt]\frac12\end{pmatrix}.
		\]

		\item 由 $L_y=(L_+-L_-)/(2i)$，
		\[
			\boxed{L_y=\frac{\hbar}{\sqrt2}
			\begin{pmatrix}0&-i&0\\i&0&-i\\0&i&0\end{pmatrix}}.
		\]
	\end{enumerate}
}


\item 
已知一粒子处于一维谐振子势场中运动，势能为$V(x)=\frac{1}{2}kx^2(k>0)$，求：
\begin{enumerate}
	%\renewcommand{\labelenumi}{\arabic{enumi}.}
	% A(\Alph) a(\alph) I(\Roman) i(\roman) 1(\arabic)
	%设定全局标号series=example	%引用全局变量resume=example
	%[topsep=-0.3em,parsep=-0.3em,itemsep=-0.3em,partopsep=-0.3em]
	%可使用leftmargin调整列表环境左边的空白长度 [leftmargin=0em]
	\item
粒子的基态本征函数$\psi_0(x)$；

\item 
若势场突然变为$V'(x)=kx^2$，则粒子仍然处于基态的概率。
\end{enumerate}

提示：用湮灭算符$\hat{a}_-=\sqrt{\frac{m\omega}{2\hbar}}(\hat{x}+\frac{i}{m\omega}\hat{p}),\sqrt2=1.414,\sqrt[4]{2}=1.189$。


\banswer{
	令 $\omega_1=\sqrt{k/m}$。由 $a_-\psi_0=0$，即
	\[
		\left(x+\frac{\hbar}{m\omega_1}\frac{d}{dx}\right)\psi_0=0,
	\]
	得到归一化基态
	\[
		\boxed{\psi_0(x)=\left(\frac{m\omega_1}{\pi\hbar}\right)^{1/4}
		\exp\!\left(-\frac{m\omega_1x^2}{2\hbar}\right)}.
	\]
	势能突变为 $kx^2=(1/2)m\omega_2^2x^2$ 后，$\omega_2=\sqrt2\,\omega_1$。突变瞬间波函数不变，先计算两个 Gaussian 基态的重叠，
	\[
	 \langle0,\omega_2|0,\omega_1\rangle
	 =\left(\frac{m^2\omega_1\omega_2}{\pi^2\hbar^2}\right)^{1/4}
	 \int_{-\infty}^\infty e^{-m(\omega_1+\omega_2)x^2/(2\hbar)}dx
	 =\frac{\sqrt2(\omega_1\omega_2)^{1/4}}{\sqrt{\omega_1+\omega_2}}.
	\]
	取模平方，处于新势场基态的概率为
	\begin{align*}
		P_0&=|\langle0,\omega_2|0,\omega_1\rangle|^2
		=\frac{2\sqrt{\omega_1\omega_2}}{\omega_1+\omega_2}\\
		&=\boxed{\frac{2^{5/4}}{1+\sqrt2}}
		\simeq0.9852.
	\end{align*}
}



\item 	
（$30'$）若已知$[\hat{a}_i,\hat{a}_j]=[\hat{a}_i^\dagger,\hat{a}_j^\dagger]=0,[\hat{a}_i,\hat{a}_j^\dagger]=\delta_{ij}$，其中$i,j=1,2$。设$J_x=\frac{1}{2}(\hat{a}^\dagger_1\hat{a}_2+\hat{a}^\dagger_2\hat{a}_1),J_y=\frac{i}{2}(\hat{a}^\dagger_1\hat{a}_2-\hat{a}^\dagger_2\hat{a}_1),J_z=\frac{1}{2}(\hat{a}^\dagger_2\hat{a}_2-\hat{a}^\dagger_1\hat{a}_1)$，求：
\begin{enumerate}
	%\renewcommand{\labelenumi}{\arabic{enumi}.}
	% A(\Alph) a(\alph) I(\Roman) i(\roman) 1(\arabic)
	%设定全局标号series=example	%引用全局变量resume=example
	%[topsep=-0.3em,parsep=-0.3em,itemsep=-0.3em,partopsep=-0.3em]
	%可使用leftmargin调整列表环境左边的空白长度 [leftmargin=0em]
	\item
$J_x,J_y,J_z$的关系式；

\item 
$J^2=J_x^2+J_y^2+J_z^2$，试用$\hat{a}_1,\hat{a}^\dagger_1,\hat{a}_2,\hat{a}_2^\dagger$表示$J^2$。
\end{enumerate}

\banswer{
令 $N_i=a_i^\dagger a_i$、$N=N_1+N_2$，并设
$K_+=a_2^\dagger a_1$、$K_-=a_1^\dagger a_2$。
题给定义等价于 $J_x=(K_++K_-)/2$、$J_y=(K_+-K_-)/(2i)$、
$J_z=(N_2-N_1)/2$；要注意题目与另一种常用约定的指标方向不同。
\begin{enumerate}
\item 不同模的算符相互对易，故
\[
 K_+K_-=N_2(N_1+1),\quad K_-K_+=N_1(N_2+1),\quad
 [K_+,K_-]=N_2-N_1=2J_z.
\]
由 $[N_i,a_j]=-\delta_{ij}a_j$、$[N_i,a_j^\dagger]=\delta_{ij}a_j^\dagger$，
还得 $[J_z,K_\pm]=\pm K_\pm$。代入线性组合，
\[
 [J_x,J_y]=\frac1{4i}[K_++K_-,K_+-K_-]=iJ_z.
\]
另外两式同理为
\[
 \boxed{[J_x,J_y]=iJ_z,\quad[J_y,J_z]=iJ_x,\quad[J_z,J_x]=iJ_y}.
\]
题中的 $J_i$ 是无量纲算符，故右端没有额外的 $\hbar$。
\item 先求平方和中前两项：
\[
 J_x^2+J_y^2=\frac12(K_+K_-+K_-K_+)
            =N_1N_2+\frac{N_1+N_2}{2}.
\]
再加 $J_z^2=(N_2-N_1)^2/4$，得到
\begin{align*}
 J^2&=N_1N_2+\frac N2+\frac{N_1^2+N_2^2-2N_1N_2}{4}\\
 &=\frac{N^2}{4}+\frac N2
 =\boxed{\frac14(a_1^\dagger a_1+a_2^\dagger a_2)
             (a_1^\dagger a_1+a_2^\dagger a_2+2)}.
\end{align*}
\end{enumerate}
}


\newpage
\item 
（$30'$）已知两种中微子的本征态为$|V_1\rangle$和$|V_2\rangle$ ，能量本征值为$E=pc+\frac{m_i^2c^4}{pc}$（其中$i=1,2$），电子中微子的本征态为$|V_e\rangle=\cos\theta|V_1\rangle+\sin\theta|V_2\rangle$，$\mu$子中微子的本征态为$|V_\mu\rangle=-\sin\theta|V_1\rangle+\cos\theta|V_2\rangle$，其中$\theta$是混合角。某体系中在$t=0$时，电子中微子处于态$|V_e\rangle$，求：
\begin{enumerate}
	%\renewcommand{\labelenumi}{\arabic{enumi}.}
	% A(\Alph) a(\alph) I(\Roman) i(\roman) 1(\arabic)
	%设定全局标号series=example	%引用全局变量resume=example
	%[topsep=-0.3em,parsep=-0.3em,itemsep=-0.3em,partopsep=-0.3em]
	%可使用leftmargin调整列表环境左边的空白长度 [leftmargin=0em]
	\item
$t$时刻中微子所处的状态；

\item 
$t$时刻电子中微子处于基态的概率。
	
\end{enumerate}


\banswer{
	\begin{enumerate}
		\item 质量本征态只积累各自的动力学相位，因此
		\[
			\boxed{|\psi(t)\rangle=\cos\theta\,e^{-iE_1t/\hbar}|V_1\rangle
			+\sin\theta\,e^{-iE_2t/\hbar}|V_2\rangle}.
		\]

		\item 再投影到电子中微子态，
		\[
			\mathcal A_{e\to e}(t)=\cos^2\theta\,e^{-iE_1t/\hbar}
			+\sin^2\theta\,e^{-iE_2t/\hbar}.
		\]
		故电子中微子的存活概率为
		\[
			\boxed{P_{e\to e}(t)=1-\sin^2(2\theta)
			\sin^2\!\left(\frac{(E_2-E_1)t}{2\hbar}\right)}.
		\]
		按题给能量公式，$E_2-E_1=(m_2^2-m_1^2)c^3/p$。
	\end{enumerate}
}


\item 
（$30'$）设在氚核中，质子和中子的作用表示成$V(r)=-V_0e^{-\frac{r}{a}}$，试用$\psi=e^{-\frac{\lambda r}{2a}}$（$\lambda$为变数）为试探波函数，以变分法求：
\begin{enumerate}
	%\renewcommand{\labelenumi}{\arabic{enumi}.}
	% A(\Alph) a(\alph) I(\Roman) i(\roman) 1(\arabic)
	%设定全局标号series=example	%引用全局变量resume=example
	%[topsep=-0.3em,parsep=-0.3em,itemsep=-0.3em,partopsep=-0.3em]
	%可使用leftmargin调整列表环境左边的空白长度 [leftmargin=0em]
	\item
基态能量的近似值；

\item 
若$V_0=32.7 \ \mathrm{Mev} ,a=2.16 \ \mathrm{fm}$，试确定$\lambda$的值。

\end{enumerate}

\banswer{
	令 $\beta=\lambda/(2a)$，归一化试探波函数为
	\[
		\psi_\lambda(r)=\left(\frac{\beta^3}{\pi}\right)^{1/2}e^{-\beta r}.
	\]
	由 $1=4\pi C^2\int_0^\infty r^2e^{-2\beta r}dr$
	可得 $C^2=\beta^3/\pi$。又因 $|\nabla\psi|^2=\beta^2|\psi|^2$，
	分部积分后 $\langle T\rangle=(\hbar^2/2\mu)\int|\nabla\psi|^2d^3r$。
	势能积分使用
	\[
	 \langle e^{-r/a}\rangle=4\beta^3\int_0^\infty r^2e^{-(2\beta+1/a)r}dr
	 =\frac{8\beta^3}{(2\beta+1/a)^3}.
	\]
	因此
	\[
		\langle T\rangle=\frac{\hbar^2\beta^2}{2\mu}
		=\frac{\hbar^2\lambda^2}{8\mu a^2},
		\qquad
		\left\langle e^{-r/a}\right\rangle
		=\frac{\lambda^3}{(\lambda+1)^3},
	\]
	其中 $\mu$ 是质子--中子体系的约化质量。因此变分能量为
	\[
		\boxed{E(\lambda)=\frac{\hbar^2\lambda^2}{8\mu a^2}
		-V_0\frac{\lambda^3}{(\lambda+1)^3}}.
	\]
	对 $\lambda>0$ 求导，
	\[
	 E'(\lambda)=\frac{\hbar^2\lambda}{4\mu a^2}
	             -\frac{3V_0\lambda^2}{(\lambda+1)^4}=0.
	\]
	非零驻点满足
	\[
		\boxed{\frac{(\lambda+1)^4}{\lambda}
		=\frac{12\mu a^2V_0}{\hbar^2}}.
	\]
	取 $\mu\simeq m_N/2=469.5\,\mathrm{MeV}/c^2$、$\hbar c=197.327\,\mathrm{MeV\,fm}$，题给参数使右端约为 $22.07$。方程有两个正根；给出能量极小值的是
	\[
		\boxed{\lambda\simeq1.326},
	\]
	相应变分基态能量为
	\[
		\boxed{E_0^{(\mathrm{var})}\simeq-2.15\,\mathrm{MeV}}.
	\]
	另一个小根约 $0.0564$ 对应局部极大值，不是基态变分解。
}



\end{enumerate}
